%% 6-3 積分するだけで解ける形（第14回：組 V6-2。ページ制限なし。偶数ページ）
%% ほかの巻・物理の本が指している形：
%%   物理 p1-10「数学 5-13 では、C を『はじめの時刻の状態』で決めました。…C を『両はしの状態』で決めます。これを境界条件といいます（数学 6-3）」
%%   物理 p1-12「傾きが一定の関数で、積分すると直線になります（6-3 で学ぶ『積分するだけで解ける形』）」（dp/dz＝−ρg、変数は高さ）
%%   物理 p1-13「右辺は一定なので √h は一定の速さで減ります（数学 6-3）」（おきかえてから積分する形 → もっと知りたい）
%%   物理 p1-14「右辺が r だけの式なので、積分するだけで解けます（6-3）」「dv/dx＝−b/m は一定なので、v は x の1次関数（6-3）」
%%   第14回の決まり：熱容量は C_th、熱抵抗は R_th
\makeatletter
% 本文の小見出し（太字の1行）  \bkvidhead{見出し}
\newcommand\bkvidhead[1]{\par\addvspace{6pt}\Needspace{6\baselineskip}\noindent{\sffamily\bfseries\color{\bkTT}#1}\par\nobreak\vspace{2pt}}
% 例題の途中式の行（「式」の札のあとにそろえる）
\newcommand\bkvidl{\par\vspace{1pt}\noindent\hangindent=2.6\zw\hangafter=0\relax}
% 大事な式のわく  \bkvidframe[はば]{中身}
\newcommand\bkvidframe[2][100mm]{\par\addvspace{4pt}{\centering\tikz\node[draw=\bkTT, line width=0.9pt, fill=bkTheme!8, rounded corners=1.5mm, inner xsep=5mm, inner ysep=2.5mm, align=center, text width=#1]{#2};\par}\addvspace{4pt}\noindent\ignorespaces}
% ルールのカード  \bkvidcard{(x,y)}{記号}{式}{説明}
\newcommand\bkvidcard[4]{\begin{scope}[shift={#1}]
  \filldraw[fill=white, draw=bkInk!60, line width=0.7pt, rounded corners=1.5mm] (0,0) rectangle (38,23);
  \node[font=\sffamily\bfseries\footnotesize, anchor=north west] at (1.2,22.4) {#2};
  \node[anchor=center, font=\small] at (19,11.6) {#3};
  \node[font=\sffamily\footnotesize, anchor=south] at (19,1) {#4};
\end{scope}}
\makeatother

\bkchapter{6}{3}{積分するだけで\\解ける形}{\bklevel{3}}{積分するだけで解ける形}
\bkmichishirube{5-11,5-12,5-13,6-1,6-2}{6-4,6-11,8-13}

\begin{goal}
\begin{itemize}[label={\color{\bkTT}●}, leftmargin=1.6\zw, itemsep=0pt, topsep=1pt]
\item \bklevel{3}\enspace 右辺に、求めたい量そのものが入っていないルールを見分けて、積分して解ける。
\item \bklevel{3}\enspace 時間でない変数（深さ・場所）のルールも、同じように積分して解ける。
\item \bklevel{3}\enspace 両はしの条件（境界条件）から、積分定数を決められる。
\end{itemize}
\end{goal}

\begin{hatena}
\noindent\begin{minipage}[c]{0.64\linewidth}\raggedright
高さ\qty{44.1}{m}のビルの屋上から、ボールを静かにはなします。6-2 の表計算なら、0.1秒ずつ進めて、地面に着く時刻をさがせます。でも、何百行も計算しなくても、\textbf{1本の式}で「何秒後に、どこにいるか」を求められないでしょうか。
\end{minipage}\hfill
\begin{minipage}[c]{0.32\linewidth}
\centering
\begin{tikzpicture}[x=0.85mm, y=0.85mm]
  \fill[bkGray!30] (0,0) rectangle (14,36);
  \draw[bkInk!70, line width=0.7pt] (0,0) rectangle (14,36);
  \foreach \yy in {4,10,16,22,28} {\foreach \xx in {2.5,8.5} {\fill[bkBlue!25] (\xx,\yy) rectangle ++(3,3.5);}}
  \draw[bkInk, line width=0.8pt] (-4,0) -- (34,0);
  \foreach \yy in {34.2,33.2,30.2,25.2,18.2,9.2} {\fill[bkOrange!45] (19,{\yy+1.8}) circle (1.5);}
  \fill[bkOrange] (19,1.6) circle (1.5);
  \draw[bkInk, line width=0.6pt, {Stealth[length=3pt]}-{Stealth[length=3pt]}] (27,0.6) -- (27,37.2);
  \node[font=\sffamily\footnotesize, anchor=west] at (28,19) {\qty{44.1}{m}};
\end{tikzpicture}
\end{minipage}
\end{hatena}

6-2 では、どんなルールでも使えるオイラー法を学びました。でも、オイラー法の答えは数の表で、少しずつずれもあります。この章では、ルールの形によっては、\textbf{積分するだけで、ぴったりの式の答えが出せる}ことを学びます。5-11〜5-13 の積分の道具が、そのまま使えます。さらに、変数が時間でないとき（水の深さ、かべの中の場所）や、「はじめ」ではなく「両はし」で条件が決まるときにも広げます。

%% ---------------- やってみよう・思い出そう ----------------
\begin{yattemiyou}[ルールを2つの仲間に分けよう]
6つのルールがあります。右辺（＝の右がわ）に何が入っているかを見て、2つの仲間に分けましょう。
\begin{bkcenter}
\begin{tikzpicture}[x=1mm, y=1mm]
  \bkvidcard{(0,26)}{ア}{$\dfrac{\dd v}{\dd t}=9.8$}{落ちるボールの速さ}
  \bkvidcard{(41,26)}{イ}{$\dfrac{\dd T}{\dd t}=-0.05(T-20)$}{冷めるお湯}
  \bkvidcard{(82,26)}{ウ}{$\dfrac{\dd V}{\dd t}=12-0.4t$}{水そうの水の量}
  \bkvidcard{(0,0)}{エ}{$\dfrac{\dd N}{\dd t}=0.2N$}{ふえる細菌の数}
  \bkvidcard{(41,0)}{オ}{$\dfrac{\dd p}{\dd z}=-9800$}{水の中の圧力（高さ $z$）}
  \bkvidcard{(82,0)}{カ}{$\dfrac{\dd T}{\dd t}=0.3$}{ヒーターで温める水}
\end{tikzpicture}
\end{bkcenter}
\begin{enumerate}[label={\bkcircled[bkGreen]{\arabic*}}, leftmargin=2.4\zw, itemsep=1pt, topsep=3pt]
\item 左辺は「何の変わる速さ」かな？　たとえばアは、速さ $v$ の変わる速さ。
\item 右辺に、その「変わる量そのもの」（アなら $v$）が入っているかどうかを見よう。
\item 「入っていない」仲間と「入っている」仲間に分けよう。
\end{enumerate}
\end{yattemiyou}

\noindent 1つずつ見ます。
\begin{itemize}[leftmargin=1.6\zw, itemsep=1pt, topsep=2pt, label={\color{\bkTT}●}]
\item ア：右辺は9.8だけ。$v$ は入っていない。\quad カ：右辺は0.3だけ。$T$ は入っていない。
\item ウ：右辺は $12-0.4t$。時刻 $t$ は入っているが、水の量 $V$ は入っていない。
\item オ：右辺は $-9800$ だけ。圧力 $p$ は入っていない（変わるのは時間ではなく高さ $z$）。
\item イ：右辺に $T$ が入っている。\quad エ：右辺に $N$ が入っている。
\end{itemize}
答えは、「入っていない」仲間がア・ウ・オ・カ、「入っている」仲間がイ・エです。この章では、「入っていない」仲間が、積分するだけで解けることを確かめます。「入っている」仲間は、6-4 で解き方を学びます。

\begin{omoidasou}
\bkrecall{5-11}{微分すると $f$ になる関数（原始関数）を求めるのが積分。$\displaystyle\int t^n\,\dd t=\frac{t^{n+1}}{n+1}+C$。$C$ は積分定数。}
\bkrecall{5-13}{$C$ は、はじめの状態（初期条件）で決まる。2回積分すると $C$ が2つ。$x=x_0+v_0t+\frac{1}{2}at^2$。}
\bkrecall{5-9}{速さのグラフの面積は、変わった量（たまった量）。}
\bkrecall{5-18}{$\displaystyle\int e^{at}\,\dd t=\frac{1}{a}e^{at}+C$。$e^0=1$、$e^{-1}=0.368$。}
\bkrecall{6-1}{ルール $\dfrac{\dd y}{\dd t}=\cdots$ は「$y$ の変わる速さ」を決める式。答えは関数。}
\bkrecall{3-6}{2つの式から2つの文字を決める（連立方程式）。}
\end{omoidasou}

%% ---------------- 数学の道具 ----------------
\bksection{doguu}{数学の道具}

\bksubsection{1}{右辺を見れば、解き方がわかる}

\leavevmode\label{bk:6-3:miwake}右辺に、変わる量そのもの（$y$）が入っていないルールは、
\[ \frac{\dd y}{\dd t}=f(t)\qquad\text{（右辺は、時刻 $t$ だけの式か、ただの数）} \]
の形をしています。「$y$ を微分すると $f(t)$ になる」と言っているので、$y$ は $f(t)$ の原始関数です（5-11）。だから、両辺を $t$ で積分すれば、
\bkvidframe[86mm]{$\displaystyle \frac{\dd y}{\dd t}=f(t)\quad\Longrightarrow\quad y=\int f(t)\,\dd t+C$}
と答えが出ます。積分定数 $C$ は、はじめの状態などの条件で決めます（5-13）。この形を、この本では「\textbf{積分するだけで解ける形}」とよびます。\bkidx{せきぶんするだけでとけるかたち}{積分するだけで解ける形}

\bkvidhead{なぜ、右辺に $y$ があるとだめなのか}
では、イの $\dfrac{\dd T}{\dd t}=-0.05\,(T-20)$ も、両辺を積分すればよいのではないでしょうか。右辺を積分するには、右辺のグラフの面積を求めなければなりません（5-9・5-10）。ところが、右辺の $T$ は、時刻とともに変わっていく「まだわからない関数」です。グラフがわからないものの面積は、求められません。

ためしに、「右辺の $T$ を、はじめの80のままにして」積分してみます。右辺は $-0.05\times(80-20)=-3$ で一定なので、
\[ T=80-3t \]
となります。これでは、20分後に $80-60=20$\,℃、30分後には $80-90=-10$\,℃になってしまいます。ほんとうは、20\,℃に近づくだけで、20\,℃より下にはなりません（6-2）。

\noindent\begin{minipage}[c]{0.52\linewidth}\centering
\begin{tikzpicture}
\begin{axis}[bkplot, width=64mm, height=46mm, xmin=0, xmax=34, ymin=-16, ymax=90, xtick={10,20,30}, ytick={-10,20,50,80}, xlabel={$t$（分）}, ylabel={$T$（℃）}, axis x line=bottom, axis y line=left]
  \addplot[bkthin, line width=0.5pt, domain=0:32] {0};
  \addplot[bkthin, dashed, domain=0:32] {20};
  \addplot[bkcurve, domain=0:32, samples=40] {20+60*exp(-0.05*x)};
  \addplot[bkRed, line width=1.3pt, domain=0:31] {80-3*x};
\end{axis}
\end{tikzpicture}
\end{minipage}\hfill
\begin{minipage}[c]{0.45\linewidth}\raggedright
赤い直線が「$T$ を80のままにした」答え、青い線がほんとうの温度。冷める速さは $T$ が下がるとおそくなるのに、赤はずっとはじめの速さで下がり続けます。右辺に $T$ が入っているときは、積分だけでは解けないのです。
\end{minipage}

\medskip
まとめると、ルールを見たら、まず右辺を見ます。
\begin{bkfigure}
\begin{tikzpicture}[x=1mm, y=1mm, nd/.style={draw=\bkTT, line width=0.9pt, rounded corners=2mm, fill=white, font=\sffamily\small, align=center, inner sep=2.5pt, minimum height=11mm}]
  \node[nd, text width=34mm] (a) at (0,0) {ルール $\dfrac{\dd y}{\dd t}=\cdots$};
  \node[nd, text width=44mm, fill=bkTheme!10] (b) at (60,10) {右辺に $y$ がない\\→ 積分するだけ（6-3）};
  \node[nd, text width=44mm] (c) at (60,-10) {右辺に $y$ がある\\→ 表計算（6-2）・6-4〜};
  \draw[bk link] (a.east) -- (b.west);
  \draw[bk link] (a.east) -- (c.west);
\end{tikzpicture}
\bkfigcaption{ルールを見たら、まず右辺を見る。}
\end{bkfigure}

\begin{seikaku}
「右辺に $y$ がない」とは、右辺が、変わる量 $y$ によらず、変数（時刻 $t$ など）だけで決まる、ということです。右辺がただの数のとき（アやカ）も、このなかまです。どんなルールでも、表計算（6-2）なら解けます。積分するだけで解ける形は、表計算を使わずに、ぴったりの式で答えが出せる、とくべつに楽な形です。
\end{seikaku}

\bksubsection{2}{落ちるボール：5-13 は微分方程式を解いていた}

\label{bk:6-3:rakka}\bkhatenaref の答えを出します。ボールの位置を、屋上から下向きに $x$〔m〕とします（下向きを＋）。空気の抵抗が小さいとき、落ちる物の加速度は $g=\qty{9.8}{m/s^2}$ で一定です（5-5）。ルールは
\[ \frac{\dd^2x}{\dd t^2}=9.8 \]
です。右辺に $x$ も $\dfrac{\dd x}{\dd t}$ も入っていないので、積分するだけで解けます。2回微分の式なので、2回積分します（5-13 と同じ手順）。
\begin{bkhenkei}
\frac{\dd v}{\dd t}=9.8 & 速度 $v=\dfrac{\dd x}{\dd t}$ で書く（ルール） \\
v=9.8t+C_1 & 1回めの積分 \\
v=9.8t & $t=0$ で $v=0$（静かにはなす）。$C_1=0$ \\
\frac{\dd x}{\dd t}=9.8t & 速度は位置の変わる速さ \\
x=4.9t^2+C_2 & 2回めの積分（$9.8\times\frac{t^2}{2}$） \\
x=4.9t^2 & $t=0$ で $x=0$（屋上）。$C_2=0$ \\
\end{bkhenkei}
地面は $x=44.1$ なので、$4.9t^2=44.1$、$t^2=9$、$t=3$\,s（$t>0$）。\textbf{3秒後に地面に着きます}。そのときの速さは $9.8\times3=\qty{29.4}{m/s}$ です。

これは、5-13 で求めた $x=x_0+v_0t+\frac{1}{2}at^2$（$x_0=0$、$v_0=0$、$a=9.8$）と同じです。5-13 で「加速度から位置を2回積分した」のは、じつは「微分方程式 $\dfrac{\dd^2x}{\dd t^2}=a$ を解いた」ことだったのです。

\bkvidhead{表計算とくらべる}
6-2 の方法（$\Delta t=0.1$\,s、5-13 の表と同じく、位置は今の速度で進める）で計算すると、3秒後の位置は\qty{42.6}{m}になり、式の\qty{44.1}{m}より少し小さくなります。式の答えは、ずれがなく、表を何百行も作らなくても、どの時刻の値もすぐに出せます。これが、積分するだけで解ける形のよさです。

\bksubsection{3}{熱：ヒーターで温める}

\label{bk:6-3:netsu}電気ケトルで水を温めます。はじめのうちは、水と部屋の温度の差が小さいので、外へ逃げる熱はわずかです。ここでは、\textbf{熱がまったく逃げない}として考えます。

\noindent\begin{minipage}[c]{0.3\linewidth}\centering
\begin{tikzpicture}[x=0.8mm, y=0.8mm]
  \filldraw[fill=bkGray!20, draw=bkInk!80, line width=0.8pt, rounded corners=2mm] (0,0) -- (22,0) -- (20,26) -- (2,26) -- cycle;
  \fill[bkBlue!30] (1.2,1.2) -- (20.8,1.2) -- (19.5,17) -- (2.5,17) -- cycle;
  \draw[bkInk!80, line width=0.8pt] (20,22) .. controls (28,22) and (28,8) .. (21,6);
  \draw[bkInk!80, line width=0.8pt] (2,24) -- (-5,28);
  \fill[bkInk!70] (-1,-2) rectangle (23,0);
  \foreach \xx in {6,11,16} {\draw[bkRed, line width=0.8pt] (\xx,2.5) -- ++(0,2.5);}
  \node[font=\sffamily\footnotesize, anchor=north] at (11,-3) {\qty{1250}{W}};
\end{tikzpicture}
\end{minipage}\hfill
\begin{minipage}[c]{0.66\linewidth}\raggedright
ヒーターの電力が $P=\qty{1250}{W}$ なら、1秒あたり\qty{1250}{J}の熱が水に入ります（1\,W は1秒あたり1\,J。4-12・5-12）。水\qty{1}{L}（\qty{1000}{g}）の温度を1\,℃上げるには、約\qty{4200}{J}の熱が要ります（水1\,g を1\,℃上げるのに約4.2\,J 要るから）。
\end{minipage}

\medskip
この「温度を1\,℃（1\,K）上げるのに要る熱」を\textbf{熱容量}（ねつようりょう）といい、$C_{\mathrm{th}}$ と書きます。単位は J/K です。\bkidx{ねつようりょう}{熱容量 $C_{\mathrm{th}}$}熱容量が大きいほど、温まりにくく、冷めにくい物です。

\bkkigocards{\bkkigocard{$C_{\mathrm{th}}$}{シー・ティーエイチ}{熱容量。温度を1\,K（1\,℃）上げるのに要る熱〔J/K〕。th は thermal（熱の）}{heat capacity}{6-3}}{\bkkigocard{\shortstack{$y(0)$\\[2pt]$y(L)$}}{ワイ・ゼロ、ワイ・エル}{$x=0$ と $x=L$ の所での $y$ の値。両はしの値の書き方}{boundary values}{6-3}}

1秒に $P$〔J〕の熱が入り、1\,K 上げるのに $C_{\mathrm{th}}$〔J〕要るので、1秒に上がる温度は $\dfrac{P}{C_{\mathrm{th}}}$〔K/s〕です。ルールは
\[ \frac{\dd T}{\dd t}=\frac{P}{C_{\mathrm{th}}}=\frac{1250}{4200}=0.2976\cdots\approx0.298\qquad\text{（℃/s）} \]
です。右辺は一定（ただの数）なので、積分するだけで解けます。はじめ20\,℃なら、
\[ T=20+0.298\,t \]
温度は、傾き約0.298の直線で上がっていきます。100\,℃になるのは、80\,℃上がったときです。まるめる前の数でわると、
\[ t=80\div\frac{1250}{4200}=80\times\frac{4200}{1250}=268.8\approx\qty{269}{s} \]
つまり $269\,\unit{s}\div\qty{60}{s/min}\approx\qty{4.5}{min}$ 後です。実際の電気ケトルでは、ケトル自身も温まり、熱も少し逃げるので、もう少し長くかかります。熱の逃げを入れると、右辺に $T$ が入ってきます（6-4 の電気ポット）。

\bksubsection{4}{電気：一定の電流でコンデンサーに電気をためる}

\label{bk:6-3:denki}6-2 の RC 回路では、たまるにつれて電流が小さくなりました。ここでは、いつも同じ電流 $I$〔A〕を流しこむ部品（定電流回路。電子回路で作れます）で、コンデンサーに電気をためます。

コンデンサーのきまり $q=CV$ と、「たまった電気の量の変わる速さは、流れこむ電流」（6-2⑧）から、
\begin{bkhenkei}
\frac{\dd q}{\dd t}=I & たまる速さは電流（一定） \\
C\frac{\dd V}{\dd t}=I & $q=CV$（$C$ は一定） \\
\frac{\dd V}{\dd t}=\frac{I}{C} & 両辺を $C$ でわる \\
\end{bkhenkei}
右辺は一定なので、積分するだけで解けます。$I=\qty{1}{mA}=\qty{0.001}{A}$、$C=\qty{100}{\micro F}=\qty{0.0001}{F}$ なら、$\dfrac{I}{C}=\dfrac{0.001}{0.0001}=10$〔V/s〕。はじめ $V=0$ なら $V=10t$ で、電圧は1秒に10\,V ずつ、まっすぐ上がっていきます。

6-2 の RC 回路（電圧が目標の10\,V にだんだん近づく曲線）と、ここ（まっすぐ上がる直線）のちがいは、ルールの右辺に $V$ が入っているかどうかのちがいです。右辺に $V$ がないと、変わる速さが一定なので、グラフは直線になります。

\bksubsection{5}{時間で変わる入力}

\label{bk:6-3:nyuuryoku}右辺は一定でなくても、時刻 $t$ だけの式なら、積分するだけで解けます。

\bkvidhead{水そうに入る水}
水が $V_0=\qty{50}{L}$ 入った水そうに、じゃ口から水を入れます。じゃ口を少しずつしめていくので、1分あたりに入る水の量（流量。5-12）は $Q=12-0.4t$〔L/min〕とだんだんへります（$0\leqq t\leqq30$）。水そうの水の量 $V$〔L〕のルールは
\[ \frac{\dd V}{\dd t}=12-0.4t \]
です（ふえる速さは、入る流量）。両辺を積分すると（5-11）、
\begin{bkhenkei}
V=12t-0.4\times\frac{t^2}{2}+C & 項ごとに積分 \\
V=12t-0.2t^2+C & $0.4\div2=0.2$ \\
V=12t-0.2t^2+50 & $t=0$ で $V=50$ なので $C=50$ \\
\end{bkhenkei}
20分後は、$V=12\times20-0.2\times400+50=240-80+50=\qty{210}{L}$ です。

\noindent\begin{minipage}[c]{0.48\linewidth}\centering
\begin{tikzpicture}
\begin{axis}[bkplot, width=52mm, height=42mm, xmin=0, xmax=33, ymin=0, ymax=14.5, xtick={10,20,30}, ytick={4,8,12}, xlabel={$t$（min）}, ylabel={$Q$（L/min）}]
  \addplot[fill=bkOrange!25, draw=none, domain=0:20] {12-0.4*x} \closedcycle;
  \addplot[bkcurve, domain=0:30] {12-0.4*x};
  \draw[bkGray, dashed, line width=0.5pt] (axis cs:20,0) -- (axis cs:20,4);
  \node[font=\sffamily\footnotesize, text=bkOrange!70!black] at (axis cs:9,4) {160};
\end{axis}
\end{tikzpicture}
\end{minipage}\hfill
\begin{minipage}[c]{0.48\linewidth}\raggedright
グラフでも確かめます。0分から20分までに入った水は、流量のグラフの下の面積です（5-9・5-12）。台形なので、$(12+4)\times20\div2=\qty{160}{L}$。はじめの50\,L とあわせて $50+160=\qty{210}{L}$。式と同じになりました。
\end{minipage}

\bkvidhead{流れこむ電流が、わかっているとき}
6-2 の RC 回路では、流れる電流は $i=0.5\,e^{-t/2}$〔mA〕でへっていくことがわかっています（6-4 で求めます）。たまる電気の量 $q$〔mC〕は、ルール $\dfrac{\dd q}{\dd t}=0.5\,e^{-t/2}$ を積分すれば出ます。$e^{at}$ の積分は $\dfrac{1}{a}e^{at}$ で（5-18）、ここでは $a=-\frac{1}{2}$ なので $\dfrac{1}{a}=-2$ です。
\begin{bkhenkei}
q=0.5\times(-2)\,e^{-t/2}+C & $\int e^{-t/2}\,\dd t=-2e^{-t/2}+C$ \\
q=-e^{-t/2}+C & $0.5\times(-2)=-1$ \\
0=-1+C,\quad C=1 & $t=0$ で $q=0$、$e^0=1$ \\
q=1-e^{-t/2} & 答え（mC） \\
\end{bkhenkei}
2秒後は $q=1-e^{-1}=1-0.368=\qty{0.632}{mC}$。コンデンサーの電圧は $V=\dfrac{q}{C}=\dfrac{0.000632}{0.0001}=\qty{6.32}{V}$ で、5-14 の式 $V=10\,(1-e^{-1})$ と同じになりました。電流がわかっていれば、たまる量は積分するだけで出せるのです。

\bksubsection{6}{変数が時間でないとき：水の深さと圧力}

\label{bk:6-3:basho}微分方程式の変数は、時間でなくてもかまいません。プールにもぐると、深くなるほど耳がぎゅっと押されます。水の圧力（水圧）が、深いほど大きいからです。

\noindent\begin{minipage}[c]{0.4\linewidth}\centering
\begin{tikzpicture}[x=1mm, y=1mm]
  \fill[bkBlue!22] (0,0) rectangle (40,34);
  \draw[bkInk!80, line width=0.8pt] (0,38) -- (0,0) -- (40,0) -- (40,38);
  \draw[bkBlue!70!black, line width=0.8pt] (0,34) -- (40,34);
  \fill[bkBlue!45] (8,16) rectangle (24,20);
  \draw[bkInk, line width=0.6pt] (8,16) rectangle (24,20);
  \draw[bk arrow, bkRed, line width=0.9pt] (17,10) -- (17,15.5);
  \draw[bk arrow, bkRed, line width=0.9pt] (17,26) -- (17,20.5);
  \node[font=\sffamily\footnotesize, anchor=west, text=bkRed!80!black] at (18,10.5) {下から押す};
  \node[font=\sffamily\footnotesize, anchor=west, text=bkRed!80!black] at (18,25.5) {上から押す};
  \draw[bk arrow, bkInk!70, line width=0.9pt] (11,18) -- (11,10);
  \fill[bkInk!70] (11,18) circle (0.6);
  \node[font=\sffamily\footnotesize, anchor=north] at (11,9.3) {重さ};
  \node[font=\sffamily\footnotesize, anchor=west] at (24.5,18) {うすい層};
  \draw[bk arrow, bkInk, line width=0.7pt] (38,4) -- (38,30) node[anchor=south, font=\sffamily\footnotesize, fill=bkBlue!22, inner sep=1pt] {$z$};
  \node[font=\sffamily\footnotesize, anchor=north west] at (0,-1) {高さ $z$ は上向きを正};
\end{tikzpicture}
\end{minipage}\hfill
\begin{minipage}[c]{0.56\linewidth}\raggedright
水の中に、厚さ $\Delta z$ のうすい層を考えます。層は、下から押される力と、上から押される力と、自分の重さでつり合っています。下の面の圧力は、上の面より、層の重さの分だけ大きくなります。
\end{minipage}

\medskip
圧力は、\qty{1}{m^2}あたりに押す力で、単位は Pa（パスカル。1\,Pa＝1\,N/m$^2$）です。面積 \qty{1}{m^2}、厚さ $\Delta z$ の層の重さは、密度 $\rho$〔kg/m$^3$〕（ロー）を使って $\rho\times\Delta z\times g$〔N〕です（体積×密度で質量、質量×$g$ で重さ）。高さを上向きに $z$ とすると、上へ $\Delta z$ 進むと圧力は $\rho g\,\Delta z$ だけへるので、$\dfrac{\Delta p}{\Delta z}=-\rho g$。層をうすくした先（5-3）が、
\bkvidframe[60mm]{$\displaystyle \frac{\dd p}{\dd z}=-\rho g$}
です。「上へ\qty{1}{m}進むと、圧力は $\rho g$ だけへる」と読みます。ここでの変数は時間 $t$ ではなく、高さ $z$ です。それでも、右辺に $p$ が入っていない（水の $\rho$ と $g$ は一定）ので、積分するだけで解けます。

水なら $\rho=\qty{1000}{kg/m^3}$ で、$\rho g=1000\times9.8=\num{9800}$〔Pa/m〕。水面の高さを $z=0$ とし、そこで大気の圧力 $p=\num{101300}$\,Pa（1013\,hPa。天気予報の「ヘクトパスカル」）とすると、
\begin{bkhenkei}
p=-9800z+C & $z$ で積分（傾きが一定なので直線） \\
p=-9800z+101300 & $z=0$ で $p=\num{101300}$、$C=\num{101300}$ \\
\end{bkhenkei}
水面から\qty{10}{m}もぐると、$z=-10$ なので $9800\times10=\num{98000}$ がたされて、$p=\num{199300}$\,Pa です。\textbf{10\,m もぐると、まわりの圧力はおよそ2倍}になります。

\begin{seikaku}
空気の中でも $\dfrac{\dd p}{\dd z}=-\rho g$ は成り立ちます。でも、空気は押しちぢめられるので、密度 $\rho$ が圧力 $p$ によって変わります。すると右辺に $p$ が入ってくるので、積分するだけでは解けません（6-4 の形になり、気圧は高さとともに指数関数でへります）。
\end{seikaku}

\bksubsection{7}{両はしで決める：かべの中の温度}

\label{bk:6-3:kyokai}冬の部屋を考えます。部屋の中は20\,℃、外は0\,℃で、そのあいだに厚さ\qty{0.2}{m}のかべがあります。かべの中の温度は、場所によってどう変わっているでしょう。

\noindent\begin{minipage}[c]{0.46\linewidth}\centering
\begin{tikzpicture}[x=0.85mm, y=0.85mm]
  \fill[bkOrange!15] (-14,0) rectangle (0,36);
  \fill[bkBlue!12] (40,0) rectangle (54,36);
  \fill[bkGray!35] (0,0) rectangle (40,36);
  \draw[bkInk, line width=0.8pt] (0,0) rectangle (40,36);
  \node[font=\sffamily\footnotesize, align=center] at (-7,30) {部屋\\20\,℃};
  \node[font=\sffamily\footnotesize, align=center] at (47,30) {外\\0\,℃};
  \draw[bkRed, line width=1.4pt] (0,28) -- (40,4);
  \fill[bkRed] (0,28) circle (1) (40,4) circle (1);
  \foreach \xx in {10,20,30} {\draw[bk arrow, bkOrange!80!black, line width=0.7pt] (\xx-4,4.5) -- (\xx+4,4.5);}
  \node[font=\sffamily\footnotesize, anchor=south] at (13,7) {熱が流れる};
  \draw[bkInk, line width=0.6pt] (0,0) -- (0,-2) (40,0) -- (40,-2);
  \node[font=\sffamily\footnotesize, anchor=north] at (0,-2.5) {$x=0$};
  \node[font=\sffamily\footnotesize, anchor=north] at (40,-2.5) {$x=0.2$};
\end{tikzpicture}
\end{minipage}\hfill
\begin{minipage}[c]{0.5\linewidth}\raggedright
部屋の側を $x=0$、外の側を $x=0.2$（m）とします。熱は、温度の高い部屋から、かべを通って外へ流れます。暖房をつけて長い時間がたち、かべの中の温度がもう変わらなくなったとき（このような状態を\textbf{定常}といいます）を考えます。
\end{minipage}

\medskip
ルールを2つのきまりから作ります。
\begin{enumerate}[label={\bkcircled[\bkTF]{\arabic*}}, leftmargin=2.4\zw, itemsep=2pt, topsep=2pt]
\item \textbf{熱は、温度の傾きに比例して流れる。}かべの中のある場所で、温度が場所とともに急に下がっているほど（傾き $\dfrac{\dd T}{\dd x}$ の大きさが大きいほど）、たくさんの熱が流れます。フーリエ（1822年）がまとめたきまりで、フーリエの法則といいます。\bkidx{ふーりえのほうそく}{フーリエの法則}
\item \textbf{定常なら、どの場所でも同じ量の熱が流れる。}もし、ある層に入る熱が出る熱より多ければ、その層はだんだん温まります。温度が変わらない（定常）ので、入る熱と出る熱は等しく、熱の流れる量はかべのどこでも同じです。
\end{enumerate}
①と②から、温度の傾き $\dfrac{\dd T}{\dd x}$ は、かべのどこでも同じです。傾きが変わらないので、傾きの変わる速さは0です。
\bkvidframe[70mm]{$\displaystyle \frac{\dd^2T}{\dd x^2}=0$}
右辺に $T$ がないので、積分するだけで解けます。2回積分します。
\begin{bkhenkei}
\frac{\dd T}{\dd x}=C_1 & 1回めの積分（0の原始関数は定数） \\
T=C_1x+C_2 & 2回めの積分。1次関数（直線） \\
\end{bkhenkei}
$C_1$ と $C_2$ を決めるのに、はじめの時刻の状態は使えません（時間の式ではないから）。わかっているのは、\textbf{かべの両はしの温度}です。
\begin{bkhenkei}
20=C_1\times0+C_2,\quad C_2=20 & 部屋の側：$x=0$ で $T=20$ \\
0=C_1\times0.2+20 & 外の側：$x=0.2$ で $T=0$ \\
C_1=-20\div0.2=-100 & $0.2\,C_1=-20$ を解く \\
T=20-100x & 答え \\
\end{bkhenkei}
かべの中の温度は、部屋の側の20\,℃から、外の側の0\,℃まで、まっすぐ下がっています。まん中（$x=0.1$）では $20-10=10$\,℃です。

5-13 では、$C$ を「はじめの時刻の状態」で決めました。かべでは、$C$ を「\textbf{両はしの状態}」で決めました。これを\textbf{境界条件}（きょうかいじょうけん）といいます。「境界」は「さかい目」のことで、ここではかべと部屋・外とのさかい目です。境界条件で答えを決める問題を、\textbf{境界値問題}といいます（はじめの状態で決める問題は、初期値問題）。\bkidx{きょうかいじょうけん}{境界条件}\bkidx{きょうかいちもんだい}{境界値問題}

\Needspace{17\baselineskip}
\begin{motto}[3]
\textbf{境界条件の問題を、表計算で解く：はさみうち}\par
6-2 の表計算は「はじめの値」から進める方法でした。両はしの条件しかわからないときは、どうすればよいでしょう。このかべで考えます。はじめの温度 $T=20$ はわかっていますが、はじめの傾き $s=\dfrac{\dd T}{\dd x}$ はわかりません。そこで、傾きを決めてためしに進め、外の側（$x=0.2$）の温度が0になるかを見ます。
\begin{bkcenter}
\begin{bktable}{colspec={Q[c,m,wd=24mm] Q[c,m,wd=40mm] Q[c,m,wd=18mm]}, rowsep=1pt, cells={font=\small}}
ためした傾き $s$ & $x=0.2$ の温度 & 判定 \\
$-50$ & $20-50\times0.2=10$ & 高すぎる \\
$-150$ & $20-150\times0.2=-10$ & 低すぎる \\
$-100$ & $20-100\times0.2=0$ & ぴったり \\
\end{bktable}
\end{bkcenter}
高すぎる傾きと低すぎる傾きで「はさみうち」にして、ちょうどよい傾きをさがします（2-16 の、表でためして答えに近づく方法）。式で解けない境界条件の問題も、この方法なら表計算で解けます（物理シリーズ『力学』の柱の曲がり方で使います）。
\end{motto}


\bksubsection{8}{両はしで決める：つり橋のケーブル}

\label{bk:6-3:hashi}つり橋は、2本の塔のあいだにケーブルをわたし、そこから橋の板（橋げた）をつり下げています。橋げたの重さは、横（水平）の長さ\qty{1}{m}あたり同じだけ、ケーブルにかかります。ケーブルの形を求めましょう。

\begin{bkfigure}
\begin{tikzpicture}[x=1mm, y=1mm]
  % 塔と橋げた
  \fill[bkGray!50] (-2,-30) rectangle (2,8) (98,-30) rectangle (102,8);
  \fill[bkInk!60] (-8,-26) rectangle (108,-24);
  \draw[bkBlue, line width=1.3pt, domain=0:100, samples=60] plot (\x, {0.004*\x*(\x-100)*2});
  \foreach \xx in {10,20,...,90} {\draw[bkInk!60, line width=0.4pt] (\xx,{0.008*\xx*(\xx-100)}) -- (\xx,-24);}
  \fill[bkInk] (0,0) circle (0.9) (100,0) circle (0.9);
  \draw[bkGray, dashed, line width=0.5pt] (-4,0) -- (104,0);
  \node[font=\sffamily\footnotesize, anchor=south] at (0,8.5) {$x=0$};
  \node[font=\sffamily\footnotesize, anchor=south] at (100,8.5) {$x=L$};
  \draw[bkInk, line width=0.6pt, {Stealth[length=3pt]}-{Stealth[length=3pt]}] (50,-0.6) -- (50,-19.4);
  \node[font=\sffamily\footnotesize, anchor=west, fill=white, inner sep=1pt] at (51.5,-7) {たるみ};
  \draw[bk arrow, bkInk, line width=0.6pt] (108,-4) -- (108,6) node[anchor=south, font=\sffamily\footnotesize] {$y$};
\end{tikzpicture}
\bkfigcaption{つり橋のケーブル（青）。塔の上を $y=0$ とし、$y$ は上向きを正にとる。たるみは見やすく大きくかいている。}
\end{bkfigure}

ケーブルの小さな部分（横のはば $\Delta x$）に働く力を考えます。ケーブルは、どこでもケーブルの向き（接線の向き）に引っぱられています。その力を横と上下に分けて考えます（4-6）。
\begin{enumerate}[label={\bkcircled[\bkTF]{\arabic*}}, leftmargin=2.4\zw, itemsep=2pt, topsep=2pt]
\item 横の力：横には橋げたの重さがかからないので、横の力はどこでも同じ大きさです。これを $H$ とします。
\item 上下の力：力の向きはケーブルの傾きと同じなので、上下の力÷横の力＝傾き $\dfrac{\dd y}{\dd x}$（4-6 の $\tan$）。上下の力は $H\times\dfrac{\dd y}{\dd x}$ です。
\item 小さな部分の右はしで上に引く力と、左はしで下に引く力の差が、その部分がささえる橋げたの重さ（横の長さ1\,m あたり $w$〔N/m〕なら $w\,\Delta x$）とつり合います。
\end{enumerate}
③を式にすると、
\[ H\times(\text{右はしの傾き})-H\times(\text{左はしの傾き})=w\,\Delta x \]
両辺を $H\,\Delta x$ でわると、$\dfrac{\text{傾きの変化}}{\Delta x}=\dfrac{w}{H}$。$\Delta x$ を小さくした先（5-3）が、
\bkvidframe[70mm]{$\displaystyle \frac{\dd^2y}{\dd x^2}=\frac{w}{H}\qquad\text{（一定）}$}
です。「傾きは、右へ行くほど同じ割合でふえる」というルールです。右辺は一定なので、積分するだけで解けます。

例として、塔と塔のあいだを $L=\qty{100}{m}$、$\dfrac{w}{H}=0.008$（1/m）とします。
\begin{bkhenkei}
\frac{\dd y}{\dd x}=0.008x+C_1 & 1回めの積分 \\
y=0.004x^2+C_1x+C_2 & 2回めの積分（$0.008\div2=0.004$） \\
\end{bkhenkei}
両はしの高さ（境界条件）で、$C_1$ と $C_2$ を決めます。
\begin{bkhenkei}
0=0+0+C_2,\quad C_2=0 & 左の塔：$x=0$ で $y=0$ \\
0=0.004\times100^2+100\,C_1 & 右の塔：$x=100$ で $y=0$ \\
C_1=-40\div100=-0.4 & $40+100\,C_1=0$ を解く \\
y=0.004x^2-0.4x & 答え \\
y=0.004x\,(x-100) & $0.004x$ でくくる（3-8） \\
\end{bkhenkei}
ケーブルの形は放物線です（3-9）。まん中（$x=50$）では
\[ y=0.004\times50\times(-50)=-10 \]
で、塔の上より\qty{10}{m}下がっています。左右の塔の高さという「両はしの状態」で、$C_1$ と $C_2$ が決まりました。これも境界条件です。

6-11 では、橋げたのない、ケーブル（くさり）だけをたらしたときの形を調べます。そのときは、重さが「横の長さ」あたりでなく「くさりの長さ」あたり同じになるので、形は放物線とは少しちがう曲線（懸垂線）になります。

\bksubsection{9}{条件はいくつ要る？}

\label{bk:6-3:jouken}ここまでの例を、ならべてみます。
\begin{bkcenter}
\begin{bktable}{colspec={Q[l,m,wd=30mm] Q[c,m,wd=26mm] Q[c,m,wd=12mm] Q[l,m,wd=40mm]}, rowsep=1pt, cells={font=\small}}
例 & ルール & 階数 & 決めるのに使った条件 \\
ヒーター & $\dfrac{\dd T}{\dd t}=0.298$ & 1階 & はじめの温度（1つ） \\
水の深さと圧力 & $\dfrac{\dd p}{\dd z}=-9800$ & 1階 & 水面の圧力（1つ） \\
落ちるボール & $\dfrac{\dd^2x}{\dd t^2}=9.8$ & 2階 & はじめの位置と速度（2つ） \\
かべの温度 & $\dfrac{\dd^2T}{\dd x^2}=0$ & 2階 & 両はしの温度（2つ） \\
つり橋 & $\dfrac{\dd^2y}{\dd x^2}=0.008$ & 2階 & 両はしの高さ（2つ） \\
\end{bktable}
\end{bkcenter}
1回積分するごとに、$C$ が1つ出てきます（5-13）。だから、\textbf{1階の式には条件が1つ、2階の式には条件が2つ}要ります。2階の式の2つの条件は、「はじめの位置と速度」（初期条件）のこともあれば、「両はしの値」（境界条件）のこともあります。片方のはしの値だけでは、$C$ は1つしか決まりません。

\begin{seikaku}
境界条件の問題では、条件の組み合わせによっては、答えがなかったり、1つに決まらなかったりすることがあります。たとえば、ばねのような $\dfrac{\dd^2y}{\dd x^2}=-k^2y$ の形の式で両はしを0にすると、とびとびの $k$ のときだけ0でない答えがあります。これは 6-6 の「もっと知りたい」（固有値問題）で調べます。
\end{seikaku}

\begin{motto}[4]
\textbf{おきかえると、積分するだけで解ける形になることがある}\par
ペットボトルの底に小さな穴をあけて水を入れると、水は穴から流れ出て、水位 $h$〔m〕が下がります。水位が高いほど、穴から出る水は速く出ます。あるペットボトルでは、ルールが
\[ \frac{\dd h}{\dd t}=-0.02\sqrt h \]
でした（$t$ は秒。$\sqrt h$ に比例するわけは物理シリーズ『力学』で）。右辺に $h$ が入っているので、このままでは積分だけでは解けません。そこで、$u=\sqrt h$ とおきかえます。合成関数の微分（5-6）で、
\begin{bkhenkei}
\frac{\dd u}{\dd t}=\frac{1}{2\sqrt h}\times\frac{\dd h}{\dd t} & $(\sqrt h)'=\dfrac{1}{2\sqrt h}$ に $h$ の微分をかける \\
\frac{\dd u}{\dd t}=\frac{1}{2\sqrt h}\times(-0.02\sqrt h) & ルールを入れる \\
\frac{\dd u}{\dd t}=-0.01 & $\sqrt h$ が約分できる \\
\end{bkhenkei}
右辺が一定になりました。$u=\sqrt h$ は、一定の速さでへっていきます。はじめ $h=0.25$（$u=0.5$）なら、$u=0.5-0.01t$ で、$u=0$（水がなくなる）のは $t=50$ 秒後です。$h=u^2=(0.5-0.01t)^2$ なので、半分の時間の25秒後には、水位は $0.25^2=0.0625$\,m（はじめの $\frac{1}{4}$）です。水がへるにつれて、出る速さもおそくなるのです。
\end{motto}

\begin{senmonsho}[大学の教科書ふうの書き方]
\begin{bkhonbun}
微分方程式 $y''=f(x)$ $(0<x<L)$ を，境界条件 $y(0)=a$，$y(L)=b$ のもとで解く問題を境界値問題という。一般解 $y=\displaystyle\iint f\,\dd x\,\dd x+C_1x+C_2$ の任意定数 $C_1$，$C_2$ は境界条件から定まる。
\end{bkhonbun}
\bkyaku{$y''$ は $y$ を $x$ で2回微分したもの。「$0<x<L$」は、かべやつり橋の中の場所。「境界条件」は、両はし $x=0$ と $x=L$ での値。$\iint f\,\dd x\,\dd x$ は「$f$ を2回積分したもの」。「任意定数」は積分定数 $C_1$、$C_2$ のこと。}
\end{senmonsho}

\begin{taisetsu}
\def\tsl#1{\makebox[6.4\zw][l]{\sffamily\bfseries\small\color{\bkTT}#1}}
\tsl{見分け方}右辺に、変わる量そのものが入っていない\\[3pt]
\tsl{}→ $\dfrac{\dd y}{\dd t}=f(t)$ なら $y=\displaystyle\int f(t)\,\dd t+C$\\[3pt]
\tsl{変数}時間でなく、高さや場所でもよい\\[3pt]
\tsl{条件の数}1階は1つ、2階は2つ（はじめの状態か、両はしの状態）
\end{taisetsu}

\begin{kotoba}
\small
\bkword{積分するだけで解ける形}{右辺に変わる量そのものが入っていないルール。両辺を積分すれば解ける。}
\bkword{熱容量}{温度を1\,K（1\,℃）上げるのに要る熱。$C_{\mathrm{th}}$〔J/K〕。}
\bkword{定常}{時間がたっても、ようすが変わらない状態。}
\bkword{境界条件}{両はし（さかい目）での値。積分定数を決めるのに使う。}
\bkword{境界値問題}{境界条件で答えを決める問題。はじめの状態で決める問題は初期値問題。}
\end{kotoba}

%% ---------------- 例題 ----------------
\bksection{reidai}{例題でたしかめよう}

\begin{reidai}[積分するだけで解ける形を見分ける]
次のルールのうち、積分するだけで解けるものをすべて選びましょう。\par
（1）$\dfrac{\dd y}{\dd t}=3t^2$\par
（2）$\dfrac{\dd y}{\dd t}=-2y$\par
（3）$\dfrac{\dd V}{\dd t}=\dfrac{6-V}{3}$\par
（4）$\dfrac{\dd p}{\dd z}=-12$\par
\kangaekata 右辺に、左辺で微分されている量（$y$、$V$、$p$）が入っているかを見る（①）。変数は時間でなくてもよい。\par
\shiki（1）右辺は $t$ だけ。$y$ はない。\par
\bkvidl（2）右辺に $y$ がある。\par
\bkvidl（3）右辺に $V$ がある。\par
\bkvidl（4）右辺はただの数。$p$ はない（変数は $z$）。\par
\kotae （1）と（4）
\end{reidai}

\begin{reidai}[落ちるボール]
高さ\qty{44.1}{m}の屋上から、ボールを静かにはなします。下向きに位置 $x$ をとると、ルールは $\dfrac{\dd^2x}{\dd t^2}=9.8$ です。$x$ を $t$ の式で表し、地面に着く時刻と、そのときの速さを求めましょう。\par
\kangaekata 右辺に $x$ がないので、2回積分する。$C_1$ ははじめの速度（0）、$C_2$ ははじめの位置（0）で決める（②）。速さは、1回めの積分の $v$ に、着く時刻を入れる。\par
\shiki $v=9.8t+C_1$。$t=0$ で $v=0$ より $C_1=0$\par
\bkvidl $x=4.9t^2+C_2$。$t=0$ で $x=0$ より $C_2=0$\par
\bkvidl $4.9t^2=44.1$、$t^2=9$、$t>0$ より $t=3$\par
\bkvidl 速さ：$v=9.8\times3=29.4$\par
\kotae $x=4.9t^2$、3秒後、\qty{29.4}{m/s}
\end{reidai}

\begingroup\tcbset{equal height group=bkvidjA}%
\noindent\begin{minipage}[t]{0.49\linewidth}
\begin{jibun}
\raggedright 次のルールのうち、積分するだけで解けるものをすべて選びましょう。\par
（1）$\dfrac{\dd T}{\dd t}=0.5$\par
（2）$\dfrac{\dd N}{\dd t}=0.3N$\par
（3）$\dfrac{\dd h}{\dd t}=2-0.1t$\par
（4）$\dfrac{\dd T}{\dd x}=-T$
\bkjans{（1）と（3）。（2）は右辺に $N$、（4）は右辺に $T$ がある。}
\end{jibun}
\end{minipage}\hfill
\begin{minipage}[t]{0.49\linewidth}
\begin{jibun}
\raggedright 高さ\qty{78.4}{m}の所から、ボールを静かにはなします。$\dfrac{\dd^2x}{\dd t^2}=9.8$（下向きが正）として、$x$ を $t$ の式で表し、地面に着く時刻と、そのときの速さを求めましょう。
\bkjans{$v=9.8t$、$x=4.9t^2$。$4.9t^2=78.4$ より $t^2=16$、$t=\qty{4}{s}$。速さ $9.8\times4=\qty{39.2}{m/s}$。}
\end{jibun}
\end{minipage}\par
\endgroup

\begin{reidai}[ヒーターで温める]
\qty{1}{L}の水（熱容量 $C_{\mathrm{th}}=\qty{4200}{J/K}$）を、電力\qty{1250}{W}のヒーターで温めます。熱は逃げないとします。はじめ20\,℃の水の温度 $T$ を時刻 $t$〔s〕の式で表し、100\,℃になる時刻を求めましょう。\par
\kangaekata 1秒に $P$〔J〕入り、1\,K に $C_{\mathrm{th}}$〔J〕要るので、$\dfrac{\dd T}{\dd t}=\dfrac{P}{C_{\mathrm{th}}}$（③）。右辺は一定なので積分するだけ。\par
\shiki $\dfrac{\dd T}{\dd t}=\dfrac{1250}{4200}\approx0.298$\par
\bkvidl $T=0.298\,t+C$。$t=0$ で $T=20$ より $C=20$\par
\bkvidl $t=80\div\dfrac{1250}{4200}=80\times\dfrac{4200}{1250}=268.8\approx269$\par
\bkvidl $269\,\unit{s}\div\qty{60}{s/min}\approx\qty{4.5}{min}$\par
\kotae $T=20+0.298\,t$、約269秒後（約4.5分後）
\end{reidai}

\begin{reidai}[一定の電流でコンデンサーに電気をためる]
静電容量 $C=\qty{100}{\micro F}$ のコンデンサーに、一定の電流 $I=\qty{1}{mA}$ を流しこみます。はじめ $V=0$ のとき、電圧 $V$ を $t$〔s〕の式で表し、5\,V になる時刻を求めましょう。\par
\kangaekata $\dfrac{\dd V}{\dd t}=\dfrac{I}{C}$（④）。単位をそろえる：\qty{1}{mA}＝\qty{0.001}{A}、\qty{100}{\micro F}＝\qty{0.0001}{F}。\par
\shiki $\dfrac{\dd V}{\dd t}=\dfrac{0.001}{0.0001}=10$\par
\bkvidl $V=10t+C$。$t=0$ で $V=0$ より $C=0$\par
\bkvidl $10t=5$ より $t=0.5$\par
\kotae $V=10t$、0.5秒後
\end{reidai}

\begingroup\tcbset{equal height group=bkvidjB}%
\noindent\begin{minipage}[t]{0.49\linewidth}
\begin{jibun}
\raggedright \qty{0.5}{L}の水（$C_{\mathrm{th}}=\qty{2100}{J/K}$）を、\qty{1000}{W}のヒーターで温めます。熱は逃げないとします。はじめ20\,℃のとき、$T$ を $t$〔s〕の式で表し、100\,℃になる時刻を求めましょう。
\bkjans{$\dfrac{\dd T}{\dd t}=\dfrac{1000}{2100}\approx0.476$。$T=20+0.476\,t$。$80\times\dfrac{2100}{1000}=\qty{168}{s}$（約2.8分）。}
\end{jibun}
\end{minipage}\hfill
\begin{minipage}[t]{0.49\linewidth}
\begin{jibun}
\raggedright $C=\qty{200}{\micro F}$ のコンデンサーに、一定の電流 $I=\qty{1}{mA}$ を流しこみます。はじめ $V=0$ のとき、$V$ を $t$〔s〕の式で表し、3\,V になる時刻を求めましょう。
\bkjans{$\dfrac{\dd V}{\dd t}=\dfrac{0.001}{0.0002}=5$。$V=5t$。$5t=3$ より \qty{0.6}{s}。}
\end{jibun}
\end{minipage}\par
\endgroup

\begin{reidai}[時間で変わる流量]
水が\qty{50}{L}入った水そうに、流量 $Q=12-0.4t$〔L/min〕で水を入れます。水の量 $V$〔L〕を $t$〔min〕の式で表し、20分後の量を求めましょう。\par
\kangaekata $\dfrac{\dd V}{\dd t}=12-0.4t$。右辺は $t$ だけなので、項ごとに積分（⑤）。$C$ ははじめの量。\par
\shiki $V=12t-0.4\times\dfrac{t^2}{2}+C=12t-0.2t^2+C$\par
\bkvidl $t=0$ で $V=50$ より $C=50$\par
\bkvidl $t=20$：$12\times20-0.2\times400+50=240-80+50=210$\par
\kotae $V=12t-0.2t^2+50$、20分後は\qty{210}{L}
\end{reidai}

\begin{reidai}[わかっている電流から、たまる量]
コンデンサーに流れこむ電流が $i=0.5\,e^{-t/2}$〔mA〕（$t$ は秒）です。はじめ $q=0$ のとき、たまった電気の量 $q$〔mC〕を $t$ の式で表し、2秒後の値を求めましょう（$e^{-1}=0.368$）。\par
\kangaekata $\dfrac{\dd q}{\dd t}=0.5\,e^{-t/2}$。$\displaystyle\int e^{at}\,\dd t=\frac{1}{a}e^{at}+C$（5-18）で $a=-\frac{1}{2}$、$\frac{1}{a}=-2$。\par
\shiki $q=0.5\times(-2)\,e^{-t/2}+C=-e^{-t/2}+C$\par
\bkvidl $t=0$ で $q=0$：$0=-1+C$ より $C=1$\par
\bkvidl $t=2$：$1-e^{-1}=1-0.368=0.632$\par
\kotae $q=1-e^{-t/2}$、2秒後は\qty{0.632}{mC}
\end{reidai}

\begingroup\tcbset{equal height group=bkvidjC}%
\noindent\begin{minipage}[t]{0.49\linewidth}
\begin{jibun}
\raggedright 水が\qty{30}{L}入った水そうに、流量 $Q=10-0.5t$〔L/min〕で水を入れます。$V$ を $t$ の式で表し、10分後の量を求めましょう。
\bkjans{$V=10t-0.25t^2+30$。$t=10$ で $100-25+30=\qty{105}{L}$。}
\end{jibun}
\end{minipage}\hfill
\begin{minipage}[t]{0.49\linewidth}
\begin{jibun}
\raggedright 流れこむ電流が $i=0.6\,e^{-t/3}$〔mA〕です。はじめ $q=0$ のとき、$q$〔mC〕を $t$ の式で表し、3秒後の値を求めましょう（$e^{-1}=0.368$）。
\bkjans{$\frac{1}{a}=-3$。$q=-1.8\,e^{-t/3}+C$、$C=1.8$。$q=1.8\,(1-e^{-t/3})$。3秒後 $1.8\times0.632\approx\qty{1.14}{mC}$。}
\end{jibun}
\end{minipage}\par
\endgroup

\begin{reidai}[深さと水圧]
海水（$\rho=\qty{1030}{kg/m^3}$）の中で、$\dfrac{\dd p}{\dd z}=-\rho g$ です（$z$ は上向きの高さ〔m〕）。海面（$z=0$）で $p=\num{101300}$\,Pa のとき、（1）$p$ を $z$ の式で表し、深さ\qty{30}{m}の水圧を求めましょう。（2）水圧が\num{303900}\,Pa（海面の3倍）になる深さを求めましょう（$g=9.8$）。\par
\kangaekata 変数は $z$。右辺は一定なので $z$ で積分（⑥）。$C$ は海面の圧力。深さ30\,m は $z=-30$。（2）は逆向き：式に $p$ を入れて、$z$ について解く（1次方程式。2-16）。\par
\shiki $\rho g=1030\times9.8=\num{10094}$\par
\bkvidl $p=-\num{10094}\,z+C$。$z=0$ で $C=\num{101300}$\par
\bkvidl （1）$z=-30$ を入れる：\par
\bkvidl \phantom{（1）}$\num{10094}\times30+\num{101300}=\num{404120}$\par
\bkvidl （2）$\num{101300}-\num{10094}\,z=\num{303900}$\par
\bkvidl \phantom{（2）}$-\num{10094}\,z=\num{202600}$\par
\bkvidl \phantom{（2）}$z=-\num{202600}\div\num{10094}\approx-20.1$\par
\kotae （1）$p=\num{101300}-\num{10094}\,z$、約\num{404000}\,Pa（海面のおよそ4倍）\par
\bkvidl（2）深さ約\qty{20}{m}
\end{reidai}

\begin{jibun}
\raggedright 湖の水（$\rho=\qty{1000}{kg/m^3}$）で、水面で $p=\num{101300}$\,Pa のとき、（1）深さ\qty{5}{m}の水圧と、（2）水圧が\num{202600}\,Pa（水面の2倍）になる深さを求めましょう（$g=9.8$）。
\bkjans{$p=\num{101300}-9800z$。（1）\nobreak $z=-5$ で $\num{101300}+\num{49000}=\num{150300}$\,Pa。（2）\nobreak $-9800z=\num{101300}$ より $z\approx-10.3$。深さ約\qty{10.3}{m}。}
\end{jibun}

\begin{reidai}[かべの中の温度（境界条件）]
厚さ\qty{0.2}{m}のかべで、部屋の側（$x=0$）が20\,℃、外の側（$x=0.2$）が0\,℃です。定常のとき $\dfrac{\dd^2T}{\dd x^2}=0$ です。（1）$T$ を $x$ の式で表し、部屋の側から\qty{5}{cm}の所の温度を求めましょう。（2）温度が10\,℃になる場所を求めましょう。\par
\kangaekata 2回積分して $T=C_1x+C_2$。両はしの温度（境界条件）で $C_1$、$C_2$ を決める（⑦）。5\,cm は0.05\,m。（2）は逆向き：式に $T=10$ を入れて、$x$ について解く。\par
\shiki $T=C_1x+C_2$\par
\bkvidl $x=0$ で $T=20$：$C_2=20$\par
\bkvidl $x=0.2$ で $T=0$：$0.2\,C_1+20=0$ より $C_1=-100$\par
\bkvidl （1）$x=0.05$：$20-100\times0.05=20-5=15$\par
\bkvidl （2）$20-100x=10$ より $100x=10$、$x=0.1$\par
\kotae （1）$T=20-100x$、15\,℃\par
\bkvidl（2）$x=\qty{0.1}{m}$（部屋の側から10\,cm）
\end{reidai}


\begin{jibun}
\raggedright 厚さ\qty{0.3}{m}のかべで、部屋の側（$x=0$）が22\,℃、外の側（$x=0.3$）が $-2$\,℃です。定常のとき、$T$ を $x$ の式で表し、（1）$x=0.1$ の温度と、（2）温度が6\,℃になる場所を求めましょう。
\bkjans{$C_2=22$、$0.3\,C_1+22=-2$ より $C_1=-80$。$T=22-80x$。（1）\nobreak $14\,\text{℃}$。（2）\nobreak $22-80x=6$ より $x=\qty{0.2}{m}$。}
\end{jibun}

\begin{reidai}[つり橋のケーブル（境界条件）]
塔と塔のあいだが\qty{100}{m}のつり橋で、ケーブルの形が $\dfrac{\dd^2y}{\dd x^2}=0.008$ にしたがいます。両方の塔の上（$x=0$ と $x=100$）で $y=0$ です。$y$ を $x$ の式で表し、$x=25$ の所の高さを求めましょう。\par
\kangaekata 2回積分して $y=0.004x^2+C_1x+C_2$。両はしの高さで $C_1$、$C_2$ を決める（⑧）。\par
\shiki $y=0.004x^2+C_1x+C_2$\par
\bkvidl $x=0$ で $y=0$：$C_2=0$\par
\bkvidl $x=100$ で $y=0$：$40+100\,C_1=0$ より $C_1=-0.4$\par
\bkvidl $y=0.004x^2-0.4x=0.004x\,(x-100)$\par
\bkvidl $x=25$：$0.004\times25\times(-75)=-7.5$\par
\kotae $y=0.004x\,(x-100)$、塔の上より\qty{7.5}{m}下
\end{reidai}

\begin{jibun}
\raggedright 塔と塔のあいだが\qty{200}{m}のつり橋で、$\dfrac{\dd^2y}{\dd x^2}=0.004$、両方の塔の上（$x=0$ と $x=200$）で $y=0$ です。$y$ を $x$ の式で表し、まん中と $x=50$ の所の高さを求めましょう。
\bkjans{$y=0.002x^2+C_1x+C_2$、$C_2=0$、$80+200\,C_1=0$ より $C_1=-0.4$。$y=0.002x\,(x-200)$。まん中 $\qty{-20}{m}$、$x=50$ で $\qty{-15}{m}$（塔の上より15\,m 下）。}
\end{jibun}

\begin{column}{熱の流れを数学にしたフーリエ}
⑦で使った「熱は温度の傾きに比例して流れる」きまりは、フランスのフーリエ（1768〜1830年）が、1822年の本『熱の解析的理論』にまとめたものです。フーリエは、かべや金属の棒の中を熱がどう伝わるかを、微分方程式で表して解こうとしました。そのとき、どんな形の温度の分布も、$\sin$ と $\cos$ の波をたし合わせて表せると考えました。これが第9巻で学ぶフーリエ級数（9-1）のはじまりです。
\end{column}

%% ---------------- まちがえやすいところ ----------------
\bksection{tsumazuki}{まちがえやすいところをチェック}

\begin{tsumazuki}{右辺に $y$ があるのに積分する}
\bkng{$\dfrac{\dd T}{\dd t}=-0.05\,(T-20)$ を積分して $T=-0.05\Bigl(\dfrac{T^2}{2}-20T\Bigr)+C$}
\bkok{右辺に $T$ があるので、積分だけでは解けない（6-2 の表計算か、6-4 の方法）}
右辺の $T$ は時刻とともに変わる関数なので、$t$ で積分するときに $T$ を「積分する文字」のようにあつかうことはできません。積分しているのは $t$ です。
\end{tsumazuki}

\begin{tsumazuki}{2回積分で、$C$ を1つしかつけない}
\bkng{$\dfrac{\dd^2T}{\dd x^2}=0$ から $T=C$}
\bkok{$\dfrac{\dd T}{\dd x}=C_1$、$T=C_1x+C_2$}
1回積分するたびに、新しい $C$ がつきます。2回積分したら $C$ は2つです。$C_1$ を積分すると $C_1x$ になることも忘れずに。
\end{tsumazuki}

\begin{tsumazuki}{片方のはしだけで $C$ を決めようとする}
\bkng{かべの部屋の側の20\,℃だけで、$T=C_1x+20$ の $C_1$ を決めようとする}
\bkok{外の側の0\,℃も使って、$C_1=-100$}
2階の式には、条件が2つ要ります（⑨）。片方のはしでは $C_2$ しか決まりません。もう片方のはしの条件で、$C_1$ を決めます。
\end{tsumazuki}

\begin{tsumazuki}{深さと高さの向きをまちがえる}
\bkng{$p=\num{101300}-9800z$ で、深さ10\,m に $z=10$ を入れる}
\bkok{$z$ は上向きなので、深さ10\,m は $z=-10$}
変数の向き（どちらを正にとるか）を、はじめに決めて変えないことが大切です（2-2）。深くなるほど圧力が大きくなるかを、答えで確かめましょう。
\end{tsumazuki}


%% ---------------- 物理・工学 ----------------
\begin{butsuri}
\begin{itemize}[leftmargin=1.2\zw, itemsep=3pt, topsep=0pt, label={\color{bkPurple}●}]
\item \textbf{断熱材の入ったかべ}\par\vspace{4pt}\noindent\begin{minipage}[c]{0.6\linewidth}\raggedright
家のかべには、熱を通しにくい断熱材が入っています。熱の通しやすさ（熱伝導率）は、コンクリートが約1.6、断熱材（グラスウールなど）が約0.04（W/(m·K)）と、40倍もちがいます。⑦と同じ考えで、2つの層では、熱の流れる量が同じになるように傾きが決まります。温度は、断熱材の中で大きく下がり（赤い線）、熱の逃げる量は、コンクリートだけのかべの10分の1ほどになります。
\end{minipage}\hfill
\begin{minipage}[c]{0.36\linewidth}\centering
\begin{tikzpicture}[x=0.66mm, y=0.75mm]
  \fill[bkGray!40] (0,0) rectangle (30,30);
  \fill[bkYellow!40] (30,0) rectangle (40,30);
  \draw[bkInk, line width=0.6pt] (0,0) rectangle (40,30) (30,0) -- (30,30);
  \draw[bkRed, line width=1.3pt] (0,27) -- (30,25.1) -- (40,2);
  \node[font=\sffamily\footnotesize, anchor=north] at (15,-0.5) {コンクリート};
  \node[font=\sffamily\footnotesize, anchor=north] at (36,-6) {断熱材};
  \node[font=\sffamily\footnotesize, anchor=north west] at (1,24.5) {20\,℃};
  \node[font=\sffamily\footnotesize, anchor=west] at (40.8,2) {0\,℃};
\end{tikzpicture}
\end{minipage}
\item \textbf{のこぎり波：一定の電流でためて、一気に放す}\par\vspace{4pt}\noindent\begin{minipage}[c]{0.6\linewidth}\raggedright
④のように一定の電流でコンデンサーに電気をためると、電圧はまっすぐ上がります。ある電圧になったら一気に放電させ、またためる、をくり返すと、のこぎりの歯のような電圧（のこぎり波）ができます。時間をはかる回路（タイマー）や、電圧を数に変える回路（AD変換器）に使われています。
\end{minipage}\hfill
\begin{minipage}[c]{0.36\linewidth}\centering
\begin{tikzpicture}[x=0.8mm, y=0.8mm]
  \draw[bk arrow, bkInk, line width=0.6pt] (0,0) -- (42,0) node[anchor=west, font=\sffamily\footnotesize] {$t$};
  \draw[bk arrow, bkInk, line width=0.6pt] (0,0) -- (0,20) node[anchor=south, font=\sffamily\footnotesize] {$V$};
  \draw[bkBlue, line width=1.2pt] (0,0) -- (12,15) -- (12.5,0) -- (24.5,15) -- (25,0) -- (37,15) -- (37.5,0);
\end{tikzpicture}
\end{minipage}
\item \textbf{はりのたわみ・管の中の流れ}\par\noindent
橋や建物のはりが重さでどれだけたわむか、細い管の中で水がどんな速さの分布で流れるかも、積分するだけで解ける形と境界条件で求められます（物理シリーズ『力学』）。
\end{itemize}
\end{butsuri}

%% ---------------- 章末問題 ----------------
\bkendproblems
\bkproblemgroup{基本}{3}

\bkproblem 次のルールのうち、積分するだけで解けるものをすべて選びましょう\mbox{（\bkreidairef{1}）}。\par
（1）$\dfrac{\dd y}{\dd t}=4t-1$\par
（2）$\dfrac{\dd T}{\dd t}=-0.1\,(T-25)$\par
（3）$\dfrac{\dd p}{\dd z}=-\num{11760}$\par
（4）$\dfrac{\dd q}{\dd t}=-\dfrac{q}{5}$
\bkpans{（1）と（3）。}

\bkproblem 高さ\qty{19.6}{m}の所からボールを静かにはなします。$\dfrac{\dd^2x}{\dd t^2}=9.8$（下向きが正）として、$x$ を $t$ の式で表し、地面に着く時刻と、そのときの速さを求めましょう\mbox{（\bkreidairef{2}）}。
\bkpans{$v=9.8t$、$x=4.9t^2$。$4.9t^2=19.6$ より $t=\qty{2}{s}$。速さ $9.8\times2=\qty{19.6}{m/s}$。}

\bkproblem \qty{2}{L}の水（$C_{\mathrm{th}}=\qty{8400}{J/K}$）を、\qty{1400}{W}のヒーターで温めます。熱は逃げないとします。はじめ15\,℃のとき、$T$ を $t$〔s〕の式で表し、85\,℃になる時刻を求めましょう\mbox{（\bkreidairef{3}）}。
\bkpans{$\dfrac{\dd T}{\dd t}=\dfrac{1400}{8400}=\dfrac{1}{6}$。$T=15+\dfrac{t}{6}$。$70\times6=\qty{420}{s}$（7分）。}

\bkproblem $C=\qty{50}{\micro F}$ のコンデンサーに、一定の電流 $I=\qty{0.2}{mA}$ を流しこみます。はじめ $V=0$ のとき、$V$ を $t$〔s〕の式で表し、6\,V になる時刻を求めましょう\mbox{（\bkreidairef{4}）}。
\bkpans{$\dfrac{\dd V}{\dd t}=\dfrac{0.0002}{0.00005}=4$。$V=4t$。$4t=6$ より \qty{1.5}{s}。}

\bkproblem 水が\qty{20}{L}入った水そうに、流量 $Q=8-0.2t$〔L/min〕で水を入れます。$V$ を $t$ の式で表し、20分後の量を求めましょう\mbox{（\bkreidairef{5}）}。
\bkpans{$V=8t-0.1t^2+20$。$t=20$ で $160-40+20=\qty{140}{L}$。}

\bkproblem コンデンサーに流れこむ電流が $i=2\,e^{-t}$〔mA〕です。はじめ $q=0$ のとき、$q$〔mC〕を $t$ の式で表し、1秒後の値を求めましょう（$e^{-1}=0.368$）\mbox{（\bkreidairef{6}）}。
\bkpans{$a=-1$、$\frac{1}{a}=-1$。$q=-2e^{-t}+C$、$C=2$。$q=2\,(1-e^{-t})$。1秒後 $2\times0.632\approx\qty{1.26}{mC}$。}

\bkproblemgroup{チャレンジ}{3}

\bkproblem 海水（$\rho=\qty{1030}{kg/m^3}$）で、海面で $p=\num{101300}$\,Pa のとき、水圧が海面の2倍（\num{202600}\,Pa）になる深さを求めましょう（$g=9.8$）\mbox{（\bkreidairef{7}）}。
\bkpans{$p=\num{101300}-\num{10094}z$。$\num{10094}\times(-z)=\num{101300}$ より、深さ約\qty{10.0}{m}（$z\approx-10.04$）。}

\bkproblem 厚さ\qty{0.25}{m}のかべで、部屋の側（$x=0$）が18\,℃、外の側（$x=0.25$）が $-2$\,℃です。定常のとき、$T$ を $x$ の式で表し、温度が0\,℃になる場所を求めましょう\mbox{（\bkreidairef{8}）}。
\bkpans{$C_2=18$、$0.25\,C_1+18=-2$ より $C_1=-80$。$T=18-80x$。$18-80x=0$ より $x=\qty{0.225}{m}$（部屋の側から22.5\,cm）。}

\bkproblem 塔と塔のあいだが\qty{120}{m}のつり橋で、$\dfrac{\dd^2y}{\dd x^2}=0.005$、両方の塔の上（$x=0$ と $x=120$）で $y=0$ です。$y$ を $x$ の式で表し、まん中のたるみを求めましょう\mbox{（\bkreidairef{9}）}。
\bkpans{$y=0.0025x^2+C_1x$、$36+120\,C_1=0$ より $C_1=-0.3$。$y=0.0025x\,(x-120)$。まん中 $x=60$ で $y=-9$。たるみ \qty{9}{m}。}

\begin{furikaeri}
\begin{checklist}
\item 右辺を見て、積分するだけで解ける形を見分け、解ける。\bkfurref{bk:6-3:miwake}
\item 深さや場所が変数のルールも、積分して解ける。\bkfurref{bk:6-3:basho}
\item 両はしの条件（境界条件）から、積分定数を決められる。\bkfurref{bk:6-3:kyokai}
\end{checklist}
\end{furikaeri}

\bkfukidashi{mina}{右辺に自分が入っていなければ、積分するだけ！　でも、冷めるお湯は右辺に $T$ がいた。次の章では、その「自分に比例してへる」形を式で解くよ。}
